% AMC 12A 2023 Problem 17 Strategic Analysis
% Generated using Strategic Problem Solving Framework v2.1

\documentclass[11pt]{article}
\usepackage{amsmath, amssymb, amsthm}
\usepackage{geometry}
\usepackage{enumitem}
\usepackage{tcolorbox}
\tcbuselibrary{skins,breakable}
\usepackage{xcolor}

% Page setup
\geometry{margin=1in}
\setlength{\parindent}{0pt}
\setlength{\parskip}{6pt}

% Custom colored boxes
\newtcolorbox{problembox}[1][]{
    colback=blue!5!white,
    colframe=blue!75!black,
    title=Problem Configuration Analysis,
    breakable,
    #1
}

\newtcolorbox{strategybox}[1][]{
    colback=pink!5!white,
    colframe=pink!75!black,
    title=Strategic Framework Application,
    breakable,
    #1
}

\newtcolorbox{tacticalbox}[1][]{
    colback=green!5!white,
    colframe=green!75!black,
    title=Tactical Selection,
    breakable,
    #1
}

\newtcolorbox{conversionbox}[1][]{
    colback=yellow!5!white,
    colframe=yellow!75!black,
    title=Conversion Techniques,
    breakable,
    #1
}

\newtcolorbox{verificationbox}[1][]{
    colback=orange!5!white,
    colframe=orange!75!black,
    title=Verification Strategy,
    breakable,
    #1
}

\newtcolorbox{roadmapbox}[1][]{
    colback=purple!5!white,
    colframe=purple!75!black,
    title=Solution Roadmap,
    breakable,
    #1
}

\newtcolorbox{competitionbox}[1][]{
    colback=red!5!white,
    colframe=red!75!black,
    title=Competition Strategy,
    breakable,
    #1
}

\newtcolorbox{pitfallbox}[1][]{
    colback=gray!5!white,
    colframe=gray!75!black,
    title=Common Pitfalls and Warnings,
    breakable,
    #1
}

\newtcolorbox{solutionbox}[1][]{
    colback=cyan!5!white,
    colframe=cyan!75!black,
    title=Detailed Solution,
    breakable,
    #1
}

\newtcolorbox{keyinsightbox}[1][]{
    colback=magenta!5!white,
    colframe=magenta!75!black,
    title=Key Insights,
    breakable,
    #1
}

\begin{document}

\title{\textbf{AMC 12A 2023 Problem 17}\\Strategic Analysis}
\author{Using Framework v2.1}
\date{}
\maketitle

\section*{Problem Statement}

Flora the frog starts at 0 on the number line and makes a sequence of jumps to the right. In any one jump, independent of previous jumps, Flora leaps a positive integer distance $m$ with probability $\frac{1}{2^m}$.

What is the probability that Flora will eventually land at 10?

\[\textbf{(A) } \frac{5}{512} \qquad \textbf{(B) } \frac{45}{1024} \qquad \textbf{(C) } \frac{127}{1024} \qquad \textbf{(D) } \frac{511}{1024} \qquad \textbf{(E) } \frac{1}{2}\]

\begin{problembox}
\section*{1. Problem Configuration Analysis}

\subsection*{A. Problem Classification}
\begin{itemize}[leftmargin=*]
    \item \textbf{Competition}: AMC 12A 2023
    \item \textbf{Problem Number}: 17
    \item \textbf{Difficulty Band}: Upper-mid range (Problem 17/25)
    \item \textbf{Expected Time}: 3-5 minutes
    \item \textbf{Point Value}: 6 points (standard AMC 12 scoring)
\end{itemize}

\subsection*{B. Mathematical Domain Classification}
\begin{itemize}[leftmargin=*]
    \item \textbf{Primary Domain}: Probability
    \item \textbf{Secondary Domain}: Combinatorics / Recursion
    \item \textbf{Tertiary Domain}: Series / Generating Functions
    \item \textbf{Cross-Domain Potential}: Very High (probability, counting, algebra, induction)
\end{itemize}

\subsection*{C. Problem Type}
\begin{itemize}[leftmargin=*]
    \item \textbf{Format}: Multiple choice (probability value)
    \item \textbf{Question Type}: ``What is the probability...''
    \item \textbf{Core Task}: Calculate probability of landing on specific position
    \item \textbf{Answer Form}: Simplified fraction
\end{itemize}

\subsection*{D. Given Information}
\begin{itemize}[leftmargin=*]
    \item Starting position: 0 on the number line
    \item Jump distribution: P(jump distance $m$) = $\frac{1}{2^m}$ for positive integer $m$
    \item Jumps are independent
    \item Flora only jumps to the right (positive direction)
    \item Target: Land on position 10
\end{itemize}

\subsection*{E. Constraints}
\begin{itemize}[leftmargin=*]
    \item Jumps must be positive integers
    \item Can only move right (no backtracking)
    \item Once past 10, cannot return
    \item Must land exactly on 10 (not just pass through)
\end{itemize}

\subsection*{F. Objective}
\begin{itemize}[leftmargin=*]
    \item \textbf{Find}: Probability of eventually landing on 10
    \item \textbf{Output}: Select correct probability from choices
\end{itemize}

\subsection*{G. Key Structural Features}
\begin{itemize}[leftmargin=*]
    \item \textbf{Probability distribution}: Sum of $\frac{1}{2^m}$ for $m \geq 1$ equals 1
    \item \textbf{Recursive structure}: Can land on 10 from any position $< 10$
    \item \textbf{Counting interpretation}: Count paths that land on 10
    \item \textbf{Symmetry}: Problem has elegant structure suggesting simple answer
    \item \textbf{Answer choices}: (E) $\frac{1}{2}$ is suspiciously simple
\end{itemize}

\subsection*{H. Strategic Orientation}
\begin{itemize}[leftmargin=*]
    \item \textbf{Primary Strategy}: Look for pattern or recursive formula
    \item \textbf{Key Insight}: Probability of landing on any $n \geq 1$ is constant
    \item \textbf{Expected Difficulty}: Moderate - requires insight or systematic calculation
    \item \textbf{Time Pressure}: Should be completed in 3-5 minutes
\end{itemize}
\end{problembox}

\begin{strategybox}
\section*{2. Strategic Framework Application}

\subsection*{A. Psychological Strategies}
\begin{itemize}[leftmargin=*]
    \item \textbf{Confidence Calibration}: 
    \begin{itemize}
        \item Upper-mid difficulty problem
        \item Multiple solution approaches available
        \item Simple answer (E) suggests elegant solution
    \end{itemize}
    \item \textbf{Pattern Recognition}:
    \begin{itemize}
        \item Answer $\frac{1}{2}$ is very clean
        \item Probability distribution sums to 1
        \item Recursive/self-similar structure
    \end{itemize}
    \item \textbf{Problem Familiarity}:
    \begin{itemize}
        \item Related to random walk problems
        \item Connection to geometric series
        \item Stars and bars counting technique
    \end{itemize}
\end{itemize}

\subsection*{B. Investigation Strategies}
\begin{itemize}[leftmargin=*]
    \item \textbf{Initial Exploration}:
    \begin{itemize}
        \item Calculate P(landing on 1): only one way, jump 1, probability $\frac{1}{2}$
        \item Calculate P(landing on 2): from 0$\to$2 directly ($\frac{1}{4}$) or 0$\to$1$\to$2 ($\frac{1}{2} \cdot \frac{1}{2} = \frac{1}{4}$), total $\frac{1}{2}$
        \item Notice pattern: P(landing on $n$) = $\frac{1}{2}$ for all $n \geq 1$
    \end{itemize}
    \item \textbf{Symmetry Observation}:
    \begin{itemize}
        \item Probability of landing on vs passing over should be related
        \item Each successful path has total probability $\frac{1}{2^{10}}$
        \item Count number of successful paths
    \end{itemize}
    \item \textbf{Multiple Approaches}:
    \begin{itemize}
        \item Method 1: Symmetry/intuitive argument
        \item Method 2: Recursive formula
        \item Method 3: Counting with stars and bars
        \item Method 4: Induction proof
        \item Method 5: Generating functions
    \end{itemize}
\end{itemize}

\subsection*{C. Argument Construction}
\begin{itemize}[leftmargin=*]
    \item \textbf{Logical Structure (Recursive Method)}:
    \begin{enumerate}
        \item Define $P(n)$ = probability of landing on $n$
        \item Base case: $P(1) = \frac{1}{2}$
        \item Recursive formula: $P(n) = \sum_{k=0}^{n-1} P(k) \cdot \frac{1}{2^{n-k}}$
        \item Show $P(n) = \frac{1}{2} P(n-1)$ by clever manipulation
        \item Conclude $P(n) = P(n-1) = \cdots = P(1) = \frac{1}{2}$
    \end{enumerate}
    \item \textbf{Key Deductions}:
    \begin{itemize}
        \item Sum $\sum_{m=1}^{\infty} \frac{1}{2^m} = 1$ (valid probability distribution)
        \item Each path to 10 has form: visit subset of $\{1,2,\ldots,9\}$, then jump to 10
        \item Number of subsets = $2^9 = 512$
        \item Each path has probability $\frac{1}{2^{10}}$
        \item Total: $\frac{512}{1024} = \frac{1}{2}$
    \end{itemize}
\end{itemize}

\subsection*{D. Cross-Domain Translation}
\begin{itemize}[leftmargin=*]
    \item \textbf{Combinatorial}: Count compositions of 10 into positive integers
    \item \textbf{Algebraic}: Recursive sequence with elegant solution
    \item \textbf{Probabilistic}: Conditional probability and independence
    \item \textbf{Analytical}: Generating function approach
\end{itemize}
\end{strategybox}

\begin{tacticalbox}
\section*{3. Tactical Methodology Selection}

\subsection*{A. Core Mathematical Tactics}
\begin{itemize}[leftmargin=*]
    \item \textbf{Pattern Recognition}: Calculate first few cases
    \item \textbf{Recursion}: Set up recursive formula for P(n)
    \item \textbf{Counting}: Use stars and bars or binary representation
    \item \textbf{Induction}: Prove P(n) = 1/2 for all $n \geq 1$
    \item \textbf{Symmetry}: Use probability distribution properties
\end{itemize}

\subsection*{B. Domain-Specific Tactics}
\begin{itemize}[leftmargin=*]
    \item \textbf{Probability}:
    \begin{itemize}
        \item Law of total probability
        \item Independence of jumps
        \item Conditioning on previous position
    \end{itemize}
    \item \textbf{Combinatorics}:
    \begin{itemize}
        \item Stars and bars for compositions
        \item Binary representation (visit or skip positions 1-9)
        \item Binomial coefficient sum: $\sum_{k=0}^{n} \binom{n}{k} = 2^n$
    \end{itemize}
    \item \textbf{Series}:
    \begin{itemize}
        \item Geometric series: $\sum_{m=1}^{\infty} \frac{1}{2^m} = 1$
        \item Telescoping or recursive simplification
    \end{itemize}
\end{itemize}

\subsection*{C. Tactical Priorities}
\begin{enumerate}[leftmargin=*]
    \item Recognize pattern by computing P(1), P(2) (90 seconds)
    \item Conjecture P(n) = 1/2 for all $n$ (30 seconds)
    \item Verify using counting method (90 seconds)
    \item Select answer (E) (30 seconds)
\end{enumerate}
\end{tacticalbox}

\begin{conversionbox}
\section*{4. Conversion Techniques Application}

\subsection*{A. Method 1: Counting Paths (Most Intuitive)}
\begin{itemize}[leftmargin=*]
    \item \textbf{Key observation}: To reach 10, Flora must visit some subset of $\{1,2,\ldots,9\}$, then jump to 10
    \item For each position $k \in \{1,2,\ldots,9\}$, Flora either:
    \begin{itemize}
        \item Lands on $k$ (visits it)
        \item Jumps over $k$ (doesn't visit it)
    \end{itemize}
    \item This gives $2^9 = 512$ possible paths
    \item \textbf{Example paths}:
    \begin{itemize}
        \item 0 $\to$ 10 (one jump of 10)
        \item 0 $\to$ 1 $\to$ 10 (jumps of 1, then 9)
        \item 0 $\to$ 5 $\to$ 10 (jumps of 5, then 5)
        \item 0 $\to$ 1 $\to$ 2 $\to$ 3 $\to$ 4 $\to$ 5 $\to$ 6 $\to$ 7 $\to$ 8 $\to$ 9 $\to$ 10
    \end{itemize}
    \item \textbf{Key insight}: Each path that lands on 10 has exactly the same total probability
    \item If path visits positions $k_1, k_2, \ldots, k_r$ before 10, the jumps are:
    \[0 \to k_1, \quad k_1 \to k_2, \quad \ldots, \quad k_r \to 10\]
    with sizes $k_1, k_2-k_1, \ldots, 10-k_r$ and total $k_1 + (k_2-k_1) + \cdots + (10-k_r) = 10$
    \item Probability of this path: $\frac{1}{2^{k_1}} \cdot \frac{1}{2^{k_2-k_1}} \cdots \frac{1}{2^{10-k_r}} = \frac{1}{2^{10}}$
    \item Every path has probability $\frac{1}{2^{10}} = \frac{1}{1024}$
    \item Total probability: $\frac{512}{1024} = \frac{1}{2}$
\end{itemize}

\subsection*{B. Method 2: Recursive Formula}
\begin{itemize}[leftmargin=*]
    \item Let $P(n)$ = probability of landing on position $n$
    \item Base case: $P(0) = 1$ (start at 0), $P(1) = \frac{1}{2}$
    \item Recursive formula: Can reach $n$ from any $k < n$ with one jump
    \[P(n) = \sum_{k=0}^{n-1} P(k) \cdot \frac{1}{2^{n-k}}\]
    \item For $n = 9$:
    \[P(9) = \sum_{k=0}^{8} P(k) \cdot \frac{1}{2^{9-k}}\]
    \item Divide by 2:
    \[\frac{P(9)}{2} = \sum_{k=0}^{8} P(k) \cdot \frac{1}{2^{10-k}}\]
    \item For $n = 10$:
    \[P(10) = \sum_{k=0}^{9} P(k) \cdot \frac{1}{2^{10-k}} = \frac{P(9)}{2} + P(9) \cdot \frac{1}{2} = P(9)\]
    \item By same reasoning: $P(10) = P(9) = P(8) = \cdots = P(2) = P(1) = \frac{1}{2}$
\end{itemize}

\subsection*{C. Method 3: Stars and Bars}
\begin{itemize}[leftmargin=*]
    \item To reach 10, need jumps that sum to 10: $j_1 + j_2 + \cdots + j_k = 10$ (all $j_i > 0$)
    \item Number of compositions of 10 into positive integers:
    \[\sum_{k=1}^{10} \binom{10-1}{k-1} = \sum_{k=1}^{10} \binom{9}{k-1} = \sum_{j=0}^{9} \binom{9}{j} = 2^9 = 512\]
    \item Each composition has probability $\frac{1}{2^{10}}$ (since jumps sum to 10)
    \item Total probability: $\frac{512}{1024} = \frac{1}{2}$
\end{itemize}

\subsection*{D. Method 4: Induction}
\begin{itemize}[leftmargin=*]
    \item \textbf{Claim}: $P(n) = \frac{1}{2}$ for all $n \geq 1$
    \item \textbf{Base case}: $P(1) = \frac{1}{2}$ (only way is jump 1 from 0)
    \item \textbf{Inductive step}: Assume $P(k) = \frac{1}{2}$ for all $1 \leq k < n$
    \item Then:
    \begin{align*}
    P(n) &= \sum_{k=0}^{n-1} P(k) \cdot \frac{1}{2^{n-k}}\\
    &= P(0) \cdot \frac{1}{2^n} + \sum_{k=1}^{n-1} \frac{1}{2} \cdot \frac{1}{2^{n-k}}\\
    &= \frac{1}{2^n} + \frac{1}{2} \sum_{k=1}^{n-1} \frac{1}{2^{n-k}}\\
    &= \frac{1}{2^n} + \frac{1}{2} \sum_{j=1}^{n-1} \frac{1}{2^j} \quad \text{(let } j = n-k \text{)}\\
    &= \frac{1}{2^n} + \frac{1}{2} \cdot \frac{\frac{1}{2}(1 - \frac{1}{2^{n-1}})}{1 - \frac{1}{2}}\\
    &= \frac{1}{2^n} + \frac{1}{2} \cdot (1 - \frac{1}{2^{n-1}})\\
    &= \frac{1}{2^n} + \frac{1}{2} - \frac{1}{2^n} = \frac{1}{2}
    \end{align*}
\end{itemize}

\subsection*{E. Method 5: Symmetry Argument}
\begin{itemize}[leftmargin=*]
    \item From position 0, Flora either:
    \begin{itemize}
        \item Lands on 10 eventually
        \item Jumps past 10 without landing on it
    \end{itemize}
    \item By symmetry, these two events are equally likely!
    \item \textbf{Why?} If Flora lands before 10, she faces the same problem scaled down
    \item The structure is self-similar, suggesting equal probabilities
    \item Therefore: P(land on 10) = P(jump past 10) = $\frac{1}{2}$
\end{itemize}
\end{conversionbox}

\begin{verificationbox}
\section*{5. Verification Strategy Design}

\subsection*{A. Verification Level}
\begin{itemize}[leftmargin=*]
    \item \textbf{Selected Level}: Level 3 - Pattern Verification
    \item \textbf{Reasoning}: Check first few cases match pattern
    \item \textbf{Time Budget}: 60 seconds
\end{itemize}

\subsection*{B. Verification Checklist}
\begin{itemize}[leftmargin=*]
    \item \textbf{Check P(1)}:
    \begin{itemize}
        \item Only way: 0 $\to$ 1 (jump 1)
        \item Probability: $\frac{1}{2}$ $\checkmark$
    \end{itemize}
    \item \textbf{Check P(2)}:
    \begin{itemize}
        \item Way 1: 0 $\to$ 2 (jump 2), probability $\frac{1}{4}$
        \item Way 2: 0 $\to$ 1 $\to$ 2 (jumps 1,1), probability $\frac{1}{2} \cdot \frac{1}{2} = \frac{1}{4}$
        \item Total: $\frac{1}{4} + \frac{1}{4} = \frac{1}{2}$ $\checkmark$
    \end{itemize}
    \item \textbf{Check P(3)}:
    \begin{itemize}
        \item 0 $\to$ 3: $\frac{1}{8}$
        \item 0 $\to$ 1 $\to$ 3: $\frac{1}{2} \cdot \frac{1}{4} = \frac{1}{8}$
        \item 0 $\to$ 2 $\to$ 3: $\frac{1}{4} \cdot \frac{1}{2} = \frac{1}{8}$
        \item 0 $\to$ 1 $\to$ 2 $\to$ 3: $\frac{1}{2} \cdot \frac{1}{2} \cdot \frac{1}{2} = \frac{1}{8}$
        \item Total: $\frac{4}{8} = \frac{1}{2}$ $\checkmark$
    \end{itemize}
    \item \textbf{Check counting}:
    \begin{itemize}
        \item For $n=10$: $2^9 = 512$ paths
        \item Each has probability $\frac{1}{1024}$
        \item Total: $\frac{512}{1024} = \frac{1}{2}$ $\checkmark$
    \end{itemize}
\end{itemize}

\subsection*{C. Alternative Verification Methods}
\begin{itemize}[leftmargin=*]
    \item \textbf{Method 1}: Verify recursive formula for several values
    \item \textbf{Method 2}: Check that probability distribution is valid (sums to 1)
    \item \textbf{Method 3}: Verify binomial sum: $\sum_{k=0}^{9} \binom{9}{k} = 2^9$
\end{itemize}

\subsection*{D. Answer Choice Validation}
\begin{itemize}[leftmargin=*]
    \item Our answer: $\frac{1}{2}$
    \item Matches choice: \textbf{(E)}
    \item Why not the others?
    \begin{itemize}
        \item (A), (B), (C), (D): All less than $\frac{1}{2}$, don't match the pattern
        \item Pattern shows P(n) = $\frac{1}{2}$ is universal for $n \geq 1$
    \end{itemize}
\end{itemize}
\end{verificationbox}

\begin{roadmapbox}
\section*{6. Solution Roadmap}

\subsection*{A. High-Level Solution Path}
\begin{enumerate}[leftmargin=*]
    \item \textbf{Compute first cases}: P(1) = $\frac{1}{2}$, P(2) = $\frac{1}{2}$
    \item \textbf{Recognize pattern}: P(n) = $\frac{1}{2}$ for all $n \geq 1$
    \item \textbf{Verify with counting}: $2^9$ paths, each with probability $\frac{1}{2^{10}}$
    \item \textbf{Calculate}: $\frac{2^9}{2^{10}} = \frac{1}{2}$
    \item \textbf{Select answer}: (E)
\end{enumerate}

\subsection*{B. Detailed Solution Steps}

\textbf{Step 1}: Calculate probability of landing on 1.
\begin{itemize}
    \item Only way: jump 1 from position 0
    \item Probability: $\frac{1}{2^1} = \frac{1}{2}$
\end{itemize}

\textbf{Step 2}: Calculate probability of landing on 2.
\begin{itemize}
    \item Direct: 0 $\to$ 2 (jump 2), probability $\frac{1}{4}$
    \item Via 1: 0 $\to$ 1 $\to$ 2, probability $\frac{1}{2} \cdot \frac{1}{2} = \frac{1}{4}$
    \item Total: $\frac{1}{4} + \frac{1}{4} = \frac{1}{2}$
\end{itemize}

\textbf{Step 3}: Recognize the pattern.
\begin{itemize}
    \item P(1) = P(2) = $\frac{1}{2}$
    \item Conjecture: P(n) = $\frac{1}{2}$ for all $n \geq 1$
\end{itemize}

\textbf{Step 4}: Verify with counting argument.
\begin{itemize}
    \item To reach 10, must choose which of $\{1,2,\ldots,9\}$ to visit
    \item Number of choices: $2^9 = 512$
    \item Each path has probability $\frac{1}{2^{10}} = \frac{1}{1024}$
    \item (Because jump sizes sum to 10)
    \item Total: $\frac{512}{1024} = \frac{1}{2}$
\end{itemize}

\textbf{Step 5}: Select the answer.
\begin{itemize}
    \item Answer is $\frac{1}{2}$, which is choice \textbf{(E)}
\end{itemize}

\subsection*{C. Key Decision Points}
\begin{itemize}[leftmargin=*]
    \item \textbf{Decision 1}: Compute first few cases to find pattern
    \item \textbf{Decision 2}: Use counting interpretation (binary choices)
    \item \textbf{Decision 3}: Recognize all paths have same probability
\end{itemize}

\subsection*{D. Time Management}
\begin{itemize}[leftmargin=*]
    \item Understanding problem: 30 seconds
    \item Computing P(1), P(2): 60 seconds
    \item Recognizing pattern: 30 seconds
    \item Counting argument: 90 seconds
    \item Verification: 30 seconds
    \item \textbf{Total}: 4 minutes
\end{itemize}
\end{roadmapbox}

\begin{competitionbox}
\section*{7. Competition Strategy Integration}

\subsection*{A. Problem Triage}
\begin{itemize}[leftmargin=*]
    \item \textbf{Difficulty Assessment}: Upper-medium (Problem 17/25)
    \item \textbf{Confidence Level}: High if you spot the pattern quickly
    \item \textbf{Decision}: Worth attempting - elegant answer suggests quick solution
\end{itemize}

\subsection*{B. Time Allocation}
\begin{itemize}[leftmargin=*]
    \item \textbf{Target Time}: 3-4 minutes
    \item \textbf{Maximum Time}: 5 minutes
    \item \textbf{Quick Strategy}: If you see P(1) = P(2) = $\frac{1}{2}$, guess (E)
\end{itemize}

\subsection*{C. Solution Format}
\begin{itemize}[leftmargin=*]
    \item Multiple choice: Select \textbf{(E)} $\frac{1}{2}$
    \item Verification: Check first 2-3 cases match pattern
    \item Bubble answer sheet carefully
\end{itemize}

\subsection*{D. Strategic Considerations}
\begin{itemize}[leftmargin=*]
    \item \textbf{Pattern Recognition}: Computing P(1) and P(2) reveals the answer
    \item \textbf{Answer Choice Clue}: (E) $\frac{1}{2}$ is suspiciously simple
    \item \textbf{Elimination}: Can't easily eliminate without computation
    \item \textbf{Time Saver}: Don't need rigorous proof in competition
\end{itemize}
\end{competitionbox}

\begin{pitfallbox}
\section*{Common Pitfalls and Warnings}

\subsection*{A. Conceptual Errors}
\begin{itemize}[leftmargin=*]
    \item \textbf{Counting only one path}: Forgetting multiple ways to reach 10
    \item \textbf{Wrong probability model}: Not using $\frac{1}{2^m}$ correctly
    \item \textbf{Stopping at finite sum}: Not considering all possible paths
    \item \textbf{Confusing land on vs pass through}: Must land exactly on 10
\end{itemize}

\subsection*{B. Computational Errors}
\begin{itemize}[leftmargin=*]
    \item \textbf{Arithmetic}: P(2) = $\frac{1}{4} + \frac{1}{4} = \frac{1}{2}$, not $\frac{1}{4}$
    \item \textbf{Binomial sum}: $\sum_{k=0}^{9} \binom{9}{k} = 2^9 = 512$, not 1024
    \item \textbf{Simplification}: $\frac{512}{1024} = \frac{1}{2}$, not other fractions
\end{itemize}

\subsection*{C. Strategic Errors}
\begin{itemize}[leftmargin=*]
    \item \textbf{Overcomplicating}: Don't need full recursion or generating functions
    \item \textbf{Missing pattern}: Not computing enough cases to see P(n) = $\frac{1}{2}$
    \item \textbf{Time management}: Should take 3-5 minutes max
\end{itemize}

\subsection*{D. Warning Signs}
\begin{itemize}[leftmargin=*]
    \item If you get answer other than $\frac{1}{2}$, recheck pattern
    \item If calculation is too complex, look for simpler approach
    \item The elegant answer (E) is a strong hint
\end{itemize}
\end{pitfallbox}

\begin{solutionbox}
\subsection*{A. Main Solution (Counting Method)}
\begin{enumerate}[leftmargin=*]
    \item \textbf{Observation}: To reach 10, Flora chooses which positions in $\{1,2,\ldots,9\}$ to land on
    
    \item \textbf{Binary choice}: For each position $k \in \{1,2,\ldots,9\}$:
    \begin{itemize}
        \item Either land on $k$ (visit it)
        \item Or jump over $k$ (don't visit it)
    \end{itemize}
    
    \item \textbf{Count paths}: $2^9 = 512$ possible paths
    
    \item \textbf{Key insight}: Each path has the same probability!
    \begin{itemize}
        \item If path visits positions $0 < k_1 < k_2 < \cdots < k_r < 10$
        \item Jump sizes: $k_1, k_2-k_1, \ldots, k_r-k_{r-1}, 10-k_r$
        \item Sum of jumps: $k_1 + (k_2-k_1) + \cdots + (10-k_r) = 10$
        \item Probability: $\frac{1}{2^{k_1}} \cdot \frac{1}{2^{k_2-k_1}} \cdots \frac{1}{2^{10-k_r}} = \frac{1}{2^{10}}$
    \end{itemize}
    
    \item \textbf{Total probability}:
    \[P(\text{land on 10}) = 512 \times \frac{1}{2^{10}} = \frac{512}{1024} = \frac{1}{2}\]
\end{enumerate}

\subsection*{B. Alternative Solution (Pattern Recognition)}
\begin{enumerate}[leftmargin=*]
    \item Calculate P(1): Only way is 0 $\to$ 1
    \[P(1) = \frac{1}{2}\]
    
    \item Calculate P(2):
    \begin{itemize}
        \item 0 $\to$ 2: probability $\frac{1}{4}$
        \item 0 $\to$ 1 $\to$ 2: probability $\frac{1}{2} \cdot \frac{1}{2} = \frac{1}{4}$
        \item Total: $P(2) = \frac{1}{2}$
    \end{itemize}
    
    \item Calculate P(3):
    \begin{itemize}
        \item 0 $\to$ 3: $\frac{1}{8}$
        \item 0 $\to$ 1 $\to$ 3: $\frac{1}{2} \cdot \frac{1}{4} = \frac{1}{8}$
        \item 0 $\to$ 2 $\to$ 3: $\frac{1}{4} \cdot \frac{1}{2} = \frac{1}{8}$
        \item 0 $\to$ 1 $\to$ 2 $\to$ 3: $\frac{1}{2} \cdot \frac{1}{2} \cdot \frac{1}{2} = \frac{1}{8}$
        \item Total: $P(3) = \frac{1}{2}$
    \end{itemize}
    
    \item Recognize pattern: $P(n) = \frac{1}{2}$ for all $n \geq 1$
    
    \item Therefore: $P(10) = \frac{1}{2}$
\end{enumerate}

\subsection*{C. Verification}
\begin{itemize}[leftmargin=*]
    \item Pattern: P(1) = P(2) = P(3) = $\frac{1}{2}$ $\checkmark$
    \item Counting: $2^9 = 512$ paths, each probability $\frac{1}{1024}$ $\checkmark$
    \item Total: $\frac{512}{1024} = \frac{1}{2}$ $\checkmark$
\end{itemize}
\end{solutionbox}

\begin{keyinsightbox}
\textbf{Key Insight}: This problem has a beautiful result: the probability of landing on ANY positive integer $n$ is always $\frac{1}{2}$!

\textbf{Why this works}:
\begin{itemize}
    \item To reach position $n$, Flora must choose which of positions $\{1, 2, \ldots, n-1\}$ to visit
    \item This is a binary choice for each position: visit or skip
    \item Number of paths: $2^{n-1}$
    \item Each path to $n$ has jump sizes that sum to $n$
    \item Therefore, each path has probability $\frac{1}{2^n}$
    \item Total probability: $\frac{2^{n-1}}{2^n} = \frac{1}{2}$
\end{itemize}

\textbf{General Pattern}: For probability distribution P(jump $m$) = $\frac{1}{2^m}$:
\begin{itemize}
    \item Distribution is valid: $\sum_{m=1}^{\infty} \frac{1}{2^m} = 1$
    \item Probability of landing on any $n \geq 1$ is $\frac{1}{2}$
    \item Probability of jumping past $n$ (never landing on it) is also $\frac{1}{2}$
    \item This creates a beautiful symmetry
\end{itemize}

\textbf{Recursive interpretation}:
\[P(n) = \sum_{k=0}^{n-1} P(k) \cdot \frac{1}{2^{n-k}}\]
can be rearranged to show $P(n) = P(n-1)$ for $n \geq 2$, confirming the constant value.

\textbf{Why answer (E) is elegant}: The simple answer $\frac{1}{2}$ reflects the deep symmetry in the problem structure. From any position, the probability of landing on the next target equals the probability of skipping it entirely.
\end{keyinsightbox}

\section*{Final Answer}

The probability that Flora will eventually land at 10 is:

\[\boxed{\textbf{(E) } \frac{1}{2}}\]

\textbf{Solution Summary}:
\begin{itemize}
    \item Each path to 10 visits some subset of $\{1,2,\ldots,9\}$
    \item Number of subsets: $2^9 = 512$
    \item Each path has probability $\frac{1}{2^{10}} = \frac{1}{1024}$
    \item Total probability: $\frac{512}{1024} = \frac{1}{2}$
\end{itemize}

\section*{Confidence Assessment}
\begin{itemize}[leftmargin=*]
    \item \textbf{Confidence Level}: 99\%
    \begin{itemize}
        \item Pattern confirmed for P(1), P(2), P(3)
        \item Counting argument is rigorous
        \item Beautiful mathematical result
        \item Multiple solution methods all agree
    \end{itemize}
    \item \textbf{Verification Applied}: Level 3 - Pattern Verification
    \begin{itemize}
        \item Verified first three cases
        \item Confirmed counting: $2^9 = 512$ paths
        \item Checked simplification: $\frac{512}{1024} = \frac{1}{2}$
    \end{itemize}
    \item \textbf{Time-Confidence Trade-off}: High confidence in 3-4 minutes
\end{itemize}

\end{document}

